% Generated draft: requires mathematical and source review.
\documentclass[11pt]{article}
\usepackage[margin=25mm]{geometry}
\usepackage{fontspec}
\setmainfont{Latin Modern Roman}
\usepackage{amsmath,amssymb,mathtools,unicode-math}
\setmathfont{Latin Modern Math}
\usepackage{enumitem,needspace,xurl,hyperref}
\setlength{\parindent}{0pt}
\setlength{\parskip}{5pt}
\setlength{\emergencystretch}{3em}
\newcommand{\N}{\mathbb{N}}
\newcommand{\Z}{\mathbb{Z}}
\newcommand{\Q}{\mathbb{Q}}
\newcommand{\R}{\mathbb{R}}
\newcommand{\C}{\mathbb{C}}
\newcommand{\F}{\mathbb{F}}
\newcommand{\A}{\mathbb{A}}
\newcommand{\PP}{\mathbb{P}}
\newcommand{\E}{\mathbb{E}}
\newcommand{\eps}{\varepsilon}
\newcommand{\dd}{\,\mathrm{d}}
\begin{document}
\section*{UnsolvedMath ID: 4200006}
\textbf{Source problem code:} catalogue entry
\subsection{Definitions and mathematical statement}
In the usual planar Birkhoff setting, a billiard table is a bounded convex domain $\Omega\subset\R^2$ with sufficient boundary regularity for the collision map $T_\Omega$ to be defined on the inward, nongrazing unit-velocity collision phase space $\mathcal P_\Omega$. The map follows a straight flight to the next boundary point and reflects specularly. The source does not state the exact smoothness or strict-convexity class of the tables in its first question.

\textbf{Conjugacy question (Victor Guillemin).} If there is a homeomorphism $S:\mathcal P_{\Omega_2}\to\mathcal P_{\Omega_1}$ such that
\[
 T_{\Omega_1}=S\circ T_{\Omega_2}\circ S^{-1},
\]
must the tables be similar, that is,
\[
 \Omega_1=aR\Omega_2+b
 \quad\text{for some }a>0,\ R\in O(2),\ b\in\R^2?
\]
Only a homeomorphism is required by the source; no differentiability, symplectic condition, length marking, or preservation of flight times is added. The conjugacy is of the collision maps rather than of a time-parameterized billiard flow.

\textbf{Related drum question.} Can one hear the shape of a real analytic convex drum? Under the customary Dirichlet-spectrum interpretation, let $\Omega_1,\Omega_2$ be bounded convex planar domains with real analytic boundaries, and let $\{\lambda_j(\Omega_i)\}_{j\geq1}$ be their Dirichlet Laplacian eigenvalues including multiplicity. Does
\[
 \lambda_j(\Omega_1)=\lambda_j(\Omega_2)\ \text{for every }j
 \quad\Longrightarrow\quad
 \Omega_1=R\Omega_2+b
\]
hold? The source does not specify the spectral boundary condition or allowable equivalence. Congruence is the natural exact-spectrum target; arbitrary scaling would rescale the spectrum. This is an explicit conventional reading, not an added claim that the excerpt selected Dirichlet data.

The supplied research summary reports both general questions still open, while describing partial rigidity or determination results in analytic, symmetric, near-integrable, or chaotic settings using spectral or marked-length information. It supplies no precise theorem or citation for those subclasses. Such additional data or restrictions are not silently appended to the general topological-conjugacy question.
\subsection{Short English statement}
Does topological conjugacy of planar billiard collision maps force their tables to be similar? Separately, does the spectrum determine a real analytic convex drum up to congruence? The table regularity and spectral boundary condition need clarification; reported special-class rigidity does not settle the general questions.
\subsection{Sources}
\begin{itemize}
\item Local UnsolvedMath registry, permanent ID 4200006, source code catalogue entry, ``Conjugacy''; proposer: Victor Guillemin. Supplied status partially solved, classification PARTIAL-PROGRESS, updated 2026-08-06.
\item Knill, \emph{Open problems in Hamiltonian dynamics} (2000), item ``6) CONJUGACY,'' \url{https://people.math.harvard.edu/~knill/seminars/intr/}. Reported HTML extraction and access: 2026-07-29. The record says the source presents an open item and no later authoritative resolution check is recorded. Rights: NEEDS\_REVIEW; public access verified but redistribution permission not established.
\item The supplied summary mentions restricted spectral/marked-length and rigidity results without identifiable papers or theorem references.
\end{itemize}
\end{document}
